\section{Conditional Gaussian distributions}
\label{sec: conditional gaussian distribution}

Theorem~\ref{thm: conditional gaussian distribution} about conditional Gaussian
distributions is well known
\citep[e.g.][Prop.~3.13.]{eatonMultivariateStatisticsVector2007}, but for your
convenience we wrote down our favorite proof of this theorem capturing
the intuition behind the statement.

\begin{theorem}[Conditional Gaussian distribution]
    \label{thm: conditional gaussian distribution}
    Let \(X\sim\normal(\mu,\Sigma)\) be a multivariate Gaussian vector where
    the covariance matrix is a block matrix of the form
    \[
        \mu = \begin{bmatrix}
            \mu_1\\ \mu_2
        \end{bmatrix}
        \quad \text{and} \quad
        \Sigma = \begin{bmatrix}
            \Sigma_{11} & \Sigma_{12}
            \\
            \Sigma_{21} & \Sigma_{22}
        \end{bmatrix},
    \]
    then the conditional distribution of \(X_2\) given \(X_1\) is
    \[
        X_2\mid X_1 \sim \normal(\mu_{2\mid 1}, \Sigma_{2\mid 1}),
    \]
    with conditional mean and variance
    \begin{align*}
        \mu_{2\mid 1} &:= \mu_2 + \Sigma_{21}\Sigma_{11}^{-1}(X_1-\mu_1)
        \\
        \Sigma_{2\mid 1} &:= \Sigma_{22} -  \Sigma_{21}\Sigma_{11}^{-1}\Sigma_{12}.
    \end{align*}
    where \(\Sigma_{11}^{-1}\) may be a generalized inverse
    \citep[cf.][]{eatonMultivariateStatisticsVector2007}.
\end{theorem}
\begin{proof}
    For the general statement we point to the standard literature
    \citep[e.g.][]{eatonMultivariateStatisticsVector2007}.
    For this proof we will assume for simplicity that \(\Sigma\) is invertible.

    Let \(\bar{X} := X-\mu\) be the centered version of \(X\). Then there exists
    a unique lower triangular matrix \(L\) such that \(\Sigma = LL^T\) (i.e.\ the Cholesky
    Decomposition). This results in the following representation\footnote{
        It is straightforward to check that \(Y := L^{-1}(X-\mu)\)
        is a multivariate Gaussian vector of centered, uncorrelated entries with variance \(1\).
        They are therefore \(\iid\) standard normal since uncorrelated multivariate Gaussian
        random variables are independent. Sometimes this representation is even
        taken as the definition of a multivariate normal distribution.
        This is the only real place we use the invertibility of \(\Sigma\) in
        its entirety, the invertibility of \(\Sigma_{11}^{-1}\) is used later on
        because we do not want to get into the business of generalized inverses
        here.
    }
    \[
        X - \mu =:
        \begin{bmatrix}
            \bar{X}_1
            \\
            \bar{X}_2
        \end{bmatrix}
        = \begin{bmatrix}
            L_{11} & 0
            \\
            L_{21} & L_{22}
        \end{bmatrix}
        \begin{bmatrix}
            Y_1 \\ Y_2
        \end{bmatrix}
        = LY
    \]
    with independent standard normal \(Y_i\), i.e. \(Y\sim\normal(0, \identity)\).
    \(L_{11}\) is invertible since \(\Sigma_{11}\) and therefore the map from
    \(Y_1\) to \(X_1\) is bijective. Conditioning on \(X_1\) is therefore
    equivalent to conditioning on \(Y_1\). But we have
    \begin{equation}
        \label{eq: underlying explicit decomposition}
        X_2 = \mu_2 + \bar{X}_2
        = \underbrace{\mu_2 + L_{21}Y_1}_{\text{conditional expectation}} + \underbrace{L_{22} Y_2}_{\text{conditional distribution}}
    \end{equation}
    So it follows that
    \[
        X_2\mid X_1 \sim  \normal(\mu_{2\mid 1}, \Sigma_{2\mid 1})
    \]
    with
    \begin{align*}
        \mu_{2\mid 1} &:= \E[X_2 \mid X_1]
        &=& \mu_2 + L_{21}Y_1
        \\
        \Sigma_{2\mid 1} &:= \Cov[X_2 \mid X_1] = \E\Bigl[ \bigl(X_2 - \E[X_2 \mid X_1]\bigr)^2 \mid X_1\Bigr]
        &=& L_{22}L_{22}^T.
    \end{align*}
    What is left to do, is to find a representation for the \(L_{ij}\)
    using the block matrices of \(\Sigma\). For this we observe
    \begin{equation}
        \label{eq: sigma to L conversion}
        \setlength\arraycolsep{4pt}
        \begin{bmatrix}
            \Sigma_{11} & \Sigma_{12}\\
            \Sigma_{21} & \Sigma_{22}
        \end{bmatrix}
        = \Sigma
        = LL^T = \begin{bmatrix}
            L_{11}L_{11}^T & L_{11}L_{21}^T
            \\
            L_{21}L_{11}^T & L_{22}L_{22}^T + L_{21}L_{21}^T
        \end{bmatrix}.
    \end{equation}
    Using \(Y_1 = L_{11}^{-1}\bar{X}_1\) and the insertion of an identity matrix this implies
    \[
        L_{21}Y_1 = L_{21}(L_{11}^T L_{11}^{-T})(L_{11}^{-1}\bar{X}_1)
        \overset{\eqref{eq: sigma to L conversion}}= \Sigma_{21}\Sigma_{11}^{-1}(X_1-\mu_1).
    \]
    resulting in the desired conditional expectation \(\mu_{2\mid 1}\). The conditional variance follows alike
    \begin{align*}
        \Sigma_{2\mid 1} = L_{22}L_{22}^T
        \overset{\eqref{eq: sigma to L conversion}}&= \Sigma_{22} - L_{21}L_{21}^T
        \\
        &= \Sigma_{22} - \underbrace{L_{21}(L_{11}^T}_{\overset{\eqref{eq: sigma to L conversion}}=\Sigma_{21}}\underbrace{L_{11}^{-T}) (L_{11}^{-1}}_{\overset{\eqref{eq: sigma to L conversion}}=\Sigma_{11}^{-1}}\underbrace{L_{11})L_{21}^T}_{\overset{\eqref{eq: sigma to L conversion}}=\Sigma_{12}}.
    \end{align*}
\end{proof}
